\ifdefined\gkver\else\def\gkver{student}\fi
\documentclass[\gkver]{gaokaozhenti}
\newcommand{\ccnum}[1]{{\fontspec{Noto Sans CJK JP}#1}}
\begin{document}

\chapter{2014年全国课标Ⅱ}
%% paperType: 理综物理部分

\begin{enumerate}
\item
%% number: 14
%% paperName: 2014年全国课标Ⅱ第 14 题 6 分
%% typeId: 1
%% type: 单选题
%% chapter: 直线运动
%% point: 运动图象
%% method: 无
%% score: 6
%% degree: 650
%% duplicateId: 0
%% body:
甲乙两汽车在一平直公路上同向行驶。在 $t=0$ 到 $t=t_{1}$ 的时间内，它们的 $v-t$ 图像如图所示。在这段时间内\xzanswer{A}
\onepicture[w=5.0cm]{figs/14.png}
\fourchoices[answer=A]
{汽车甲的平均速度比乙大}
{汽车乙的平均速度等于 $\frac{v_{1}+v_{2}}{2}$}
{甲乙两汽车的位移相同}
{汽车甲的加速度大小逐渐减小，汽车乙的加速度大小逐渐增大}
%% answer: A
%% memo:
\memoanswer{
由于图线与坐标轴所夹的面积表示物体的位移，在 $0-t_{1}$ 时间内，甲车的位移大于乙车，由 $\overline{v}=\frac{x}{t}$ 可知，甲车的平均速度大于乙车，A 正确，C 错误；因为乙车做变减速运动故平均速度不等于 $\frac{v_{1}+v_{2}}{2}$，B 错误；又图线的切线的斜率等于物体的加速度，则甲乙两车的加速度均逐渐减小，选项 D 错误。
}

\item
%% number: 15
%% paperName: 2014年全国课标Ⅱ第 15 题 6 分
%% typeId: 1
%% type: 单选题
%% chapter: 曲线运动
%% point: 平抛运动
%% method: 无
%% score: 6
%% degree: 450
%% duplicateId: 0
%% body:
取水平地面为重力势能零点。一物块从某一高度水平抛出，在抛出点其动能与重力势能恰好相等。不计空气阻力，该物块落地时的速度方向与水平方向的夹角为\xzanswer{B}
\fourchoices[answer=B]
{$\frac{\pi}{6}$}
{$\frac{\pi}{4}$}
{$\frac{\pi}{3}$}
{$\frac{5\pi}{12}$}
%% answer: B
%% memo:
\memoanswer{
设物体水平抛出的初速度为 $v_{0}$，抛出时的高度为 $h$，则 $\frac{1}{2}mv_{0}^{2}=mgh$，故 $v_{0}=\sqrt{2gh}$。物体落地的竖直速度 $v_{y}=\sqrt{2gh}$，则落地时速度方向与水平方向的夹角 $\tan\theta=\frac{v_{y}}{v_{x}}=\frac{v_{y}}{v_{0}}=1$，则 $\theta=\frac{\pi}{4}$，选项 B 正确。
}

\item
%% number: 16
%% paperName: 2014年全国课标Ⅱ第 16 题 6 分
%% typeId: 1
%% type: 单选题
%% chapter: 功和能
%% point: 动能定理
%% method: 无
%% score: 6
%% degree: 450
%% duplicateId: 0
%% body:
一物体静止在粗糙水平地面上，现用一大小为 $F_{1}$ 的水平拉力拉动物体，经过一段时间后其速度变为 $v$，若将水平拉力的大小改为 $F_{2}$，物体从静止开始经过同样的时间后速度变为 $2v$，对于上述两个过程，用 $W_{F1}$、$W_{F2}$ 分别表示拉力 $F_{1}$、$F_{2}$ 所做的功，$W_{f1}$、$W_{f2}$ 分别表示前后两次克服摩擦力所做的功，则\xzanswer{C}
\fourchoices[answer=C]
{$W_{F2}>4W_{F1}$，$W_{f2}>2W_{f1}$}
{$W_{F2}>4W_{F1}$，$W_{f2}=2W_{f1}$}
{$W_{F2}<4W_{F1}$，$W_{f2}=2W_{f1}$}
{$W_{F2}<4W_{F1}$，$W_{f2}<2W_{f1}$}
%% answer: C
%% memo:
\memoanswer{
由于物体两次受恒力作用做匀加速运动，由于时间相等，末速度之比为 1:2，则加速度之比为 1:2，位移之比为 1:2。而摩擦力不变，由 $W_{f}=-F_{f}\cdot x$ 得：$W_{f2}=2W_{f1}$；由动能定理：$W_{F1}+W_{f1}=\frac{1}{2}mv^{2}-0$，$W_{F2}+W_{f2}=\frac{1}{2}m(2v)^{2}-0$，整理得：$4W_{F1}=W_{F2}-2W_{f1}=W_{F2}+2F_{f}\cdot x_{1}$，故 $W_{F2}<4W_{F1}$。C 正确。
}

\item
%% number: 17
%% paperName: 2014年全国课标Ⅱ第 17 题 6 分
%% typeId: 1
%% type: 单选题
%% chapter: 曲线运动
%% point: 竖直平面圆周运动
%% method: 整体与隔离
%% score: 6
%% degree: 450
%% duplicateId: 0
%% body:
如图，一质量为 $M$ 的光滑大圆环，用一细轻杆固定在竖直平面内；套在大环上质量为 $m$ 的小环（可视为质点），从大环的最高处由静止滑下。重力加速度大小为 $g$，当小环滑到大环的最低点时，大环对轻杆拉力的大小为\xzanswer{C}
\onepicture[w=3.0cm]{figs/17.png}
\fourchoices[answer=C]
{$Mg-5Mg$}
{$Mg+mg$}
{$Mg+5mg$}
{$Mg+10mg$}
%% answer: C
%% memo:
\memoanswer{
小圆环到达大圆环低端过程机械能守恒：$mg\cdot 2R=\frac{1}{2}mv^{2}$，小圆环在最低点，根据牛顿第二定律可得：$F_{N}-mg=m\frac{v^{2}}{R}$；对大圆环，由平衡可得：$F_{T}=Mg+F_{N}$，解得 $F_{T}=Mg+5mg$，选项 C 正确。
}

\item
%% number: 18
%% paperName: 2014年全国课标Ⅱ第 18 题 6 分
%% typeId: 1
%% type: 单选题
%% chapter: 曲线运动
%% point: 天体运动
%% method: 无
%% score: 6
%% degree: 450
%% duplicateId: 0
%% body:
假设地球可视为质量均匀分布的球体。已知地球表面重力加速度在两极的大小为 $g_{0}$，在赤道的大小为 $g$，地球自转的周期为 $T$，引力常量为 $G$。地球的密度为\xzanswer{B}
\fourchoices[answer=B]
{$\frac{3\pi}{GT^{2}}\frac{g_{0}-g}{g_{0}}$}
{$\frac{3\pi}{GT^{2}}\frac{g_{0}}{g_{0}-g}$}
{$\frac{3\pi}{GT^{2}}$}
{$\frac{3\pi}{GT^{2}}\frac{g_{0}}{g}$}
%% answer: B
%% memo:
\memoanswer{
地球表面的物体在两极万有引力等于重力：$G\frac{Mm}{R^{2}}=mg_{0}$，在地球的赤道上万有引力为重力和物体随地球自转所需向心力的合力：$G\frac{Mm}{R^{2}}=m\left(\frac{2\pi}{T}\right)^{2}R+mg$，地球的质量：$M=\rho'\frac{4}{3}\pi R^{3}$，联立可得：$\rho=\frac{3\pi}{GT^{2}}\frac{g_{0}}{g_{0}-g}$，选项 B 正确。
}

\item
%% number: 19
%% paperName: 2014年全国课标Ⅱ第 19 题 6 分
%% typeId: 2
%% type: 多选题
%% chapter: 电场
%% point: 电势能电势差电场力做功
%% method: 无
%% score: 6
%% degree: 650
%% duplicateId: 0
%% body:
关于静电场的电场强度和电势，下列说法正确的是\xzanswer{AD}
\fourchoices[answer=AD]
{电场强度的方向处处与等电势面垂直}
{电场强度为零的地方，电势也为零}
{随着电场强度的大小逐渐减小，电势也逐渐降低}
{任一点的电场强度总是指向该点电势降落最快的方向}
%% answer: AD
%% memo:
\memoanswer{
由于电场线与等势面垂直，而电场强度的方向为电场线的方向，故电场强度的方向与等势面垂直，A 正确；场强为零的地方电势不一定为零，如等量同种正电荷连线的中点处的场强为零，但电势大于零，B 错误；场强大小与电场线的疏密有关，而沿着电场线的方向电势逐点降低，电场线的疏密与电场线的方向没有必然联系，C 错误；任一点的电场强度方向总是和电场线方向一致，沿电场线的方向是电势降落最快的方向，D 正确。
}

\item
%% number: 20
%% paperName: 2014年全国课标Ⅱ第 20 题 6 分
%% typeId: 2
%% type: 多选题
%% chapter: 磁场
%% point: 带电粒子在磁场中的运动
%% method: 无
%% score: 6
%% degree: 650
%% duplicateId: 0
%% body:
如图所示为某磁谱仪部分构件的示意图。图中，永磁铁提供匀强磁场，硅微条径迹探测器可以探测粒子在其中运动的轨迹。宇宙射线中有大量的电子、正电子和质子。当这些粒子从上部垂直进入磁场时，下列说法正确的是\xzanswer{AC}
\onepicture[w=5.0cm]{figs/20.png}
\fourchoices[answer=AC]
{电子与正电子的偏转方向一定不同}
{电子与正电子在磁场中运动轨迹的半径一定相同}
{仅依据粒子运动轨迹无法判断该粒子是质子还是正电子}
{粒子的动能越大，它在磁场中运动轨迹的半径越小}
%% answer: AC
%% memo:
\memoanswer{
由于电子和正电子带电性相反，若入射速度方向相同时，受力方向相反，则偏转方向一定相反，A 正确；由于电子和正电子的入射速度大小未知，由 $qvB=m\frac{v^{2}}{r}$ 得 $r=\frac{mv}{qB}$ 可知，运动半径不一定相同，B 错误；虽然质子和正电子带电量及电性相同，但是两者的动量大小未知，无法判断粒子是质子还是正电子，C 正确；由 $E_{k}=\frac{1}{2}mv^{2}$，则 $r=\frac{mv}{qB}=\frac{\sqrt{2mE_{k}}}{qB}$，粒子的动能越大，它在磁场中运动轨迹的半径越大，D 错误。
}

\item
%% number: 21
%% paperName: 2014年全国课标Ⅱ第 21 题 6 分
%% typeId: 2
%% type: 多选题
%% chapter: 交流电
%% point: 变压器
%% method: 无
%% score: 6
%% degree: 450
%% duplicateId: 0
%% body:
如图，一理想变压器原、副线圈的匝数分别为 $n_{1}$、$n_{2}$。原线圈通过一理想电流表 A 接正弦交流电源，一个二极管和阻值为 $R$ 的负载电阻串联后接到副线圈的两端。假设该二极管的正向电阻为零，反向电阻为无穷大。用交流电压表测得 $a$、$b$ 端和 $c$、$d$ 端的电压分别为 $U_{ab}$ 和 $U_{cd}$，则\xzanswer{BD}
\onepicture[w=5.5cm]{figs/21.png}
\fourchoices[answer=BD]
{$U_{ab}:U_{cd}=n_{1}:n_{2}$}
{增大负载电阻的阻值 $R$，电流表的读数变小}
{负载电阻的阻值越小，$cd$ 间的电压 $U_{cd}$ 越大}
{将二极管短路，电流表的读数加倍}
%% answer: BD
%% memo:
\memoanswer{
若变压器初级电压为 $U_{ab}$，则次级电压为 $U_{2}=\frac{n_{2}}{n_{1}}U_{ab}$；由于二极管的单向导电性使得 $cd$ 间电压为 $U_{cd}=\frac{U_{2}}{2}$，故 $\frac{U_{ab}}{U_{cd}}=\frac{2n_{1}}{n_{2}}$，A 错误；增大负载的阻值 $R$，输出功率减小，则输入电流减小，即电流表的读数减小，B 正确；$cd$ 间的电压由变压器的初级电压决定，与负载电阻 $R$ 的大小无关，C 错误；若二极管短路则 $U_{cd}=U_{2}$，则次级电流会加倍，则初级电流也加倍，D 正确。
}

\item
%% number: 22
%% paperName: 2014年全国课标Ⅱ第 22 题 6 分
%% typeId: 4
%% type: 实验题
%% chapter: 电路
%% point: 电路实验
%% method: 无
%% score: 6
%% degree: 650
%% duplicateId: 0
%% body:
在伏安法测电阻的实验中，待测电阻 $R_{x}$ 约为 $200\UO$，电压表 V 的内阻为 $2\UkO$，电流表 A 的内阻约为 $10\UO$，测量电路中电流表的连接方式如图（a）或图（b）所示，结果由公式 $R_{x}=\frac{U}{I}$ 计算得出，式中 $U$ 与 $I$ 分别为电压表和电流表的示数。若将图（a）和图（b）中电路图测得的电阻值分别记为 $R_{x1}$ 和 $R_{x2}$，则\tkanswer[1.2cm]{$R_{x1}$}（填“$R_{x1}$”或“$R_{x2}$”）更接近待测电阻的真实值，且测量值 $R_{x1}$\tkanswer[1.2cm]{大于}（填“大于”“等于”或“小于”）真实值，测量值 $R_{x2}$\tkanswer[1.2cm]{小于}（填“大于”“等于”或“小于”）真实值。
\onepicture[align=right, w=4.0cm]{figs/22.png}
%% answer:
%% memo:
\memoanswer{
由于 $\frac{R_{V}}{R_{x}}=10$，$\frac{R_{x}}{R_{A}}=20$，可认为是大电阻，采取电流表内接法测量更准确，即图（a）电路测量，$R_{x1}$ 更接近待测电阻真实值；图（a）电路测量由于电流表的分压使测量值大于真实值；图（b）电路测量由于电压表的分流使测量值小于真实值。
}

\item
%% number: 23
%% paperName: 2014年全国课标Ⅱ第 23 题 9 分
%% typeId: 4
%% type: 实验题
%% chapter: 力物体的平衡
%% point: 胡克定律
%% method: 无
%% score: 9
%% degree: 550
%% duplicateId: 0
%% body:
某实验小组探究弹簧的劲度系数 $k$ 与其长度（圈数）的关系。实验装置如图（a）所示；一均匀长弹簧竖直悬挂，7 个指针 $P_{0}$、$P_{1}$、$P_{2}$、$P_{3}$、$P_{4}$、$P_{5}$、$P_{6}$ 分别固定在弹簧上距悬点 0、10、20、30、40、50、60 圈处；通过旁边竖直放置的刻度尺，可以读出指针的位置，$P_{0}$ 指向 0 刻度。设弹簧下端未挂重物时，各指针的位置记为 $x_{0}$；挂有质量为 $0.100\Ukg$ 的砝码时，各指针的位置记为 $x$。测量结果及部分计算结果如下表所示（$n$ 为弹簧的圈数，取重力加速度为 $9.80\Umsq$）。已知实验所用弹簧总圈数为 60，整个弹簧的自由长度为 $11.88\Ucm$。

\begin{center}
\begin{tabular}{|c|c|c|c|c|c|c|}
\hline
	 & $P_{1}$ & $P_{2}$ & $P_{3}$ & $P_{4}$ & $P_{5}$ & $P_{6}$ \\ \hline
	$x_{0}$（$\Ucm$） & 2.04 & 4.06 & 6.06 & 8.05 & 10.03 & 12.01 \\ \hline
	$x$（$\Ucm$） & 2.64 & 5.26 & 7.81 & 10.30 & 12.93 & 15.41 \\ \hline
	$n$ & 10 & 20 & 30 & 40 & 50 & 60 \\ \hline
	$k$（$\UNm$） & 163 & \ccnum{①} & 56.0 & 43.6 & 33.8 & 28.8 \\ \hline
	$1/k$（$\UNmn$） & 0.0061 & \ccnum{②} & 0.0179 & 0.0229 & 0.0296 & 0.0347 \\ \hline
\end{tabular}
\end{center}

\begin{enumerate}
	\item
	将表中数据补充完整：\ccnum{①}\_\_\_\_\_\_\_\_，\ccnum{②}\_\_\_\_\_\_\_\_。
	\item
	以 $n$ 为横坐标，$1/k$ 为纵坐标，在图（b）给出的坐标纸上画出 $1/k-n$ 图像。
	\item
	图（b）中画出的直线可以近似认为通过原点，若从试验中所用的弹簧截取圈数为 $n$ 的一段弹簧，该弹簧的劲度系数 $k$ 与其圈数 $n$ 的关系表达式为 $k=$\_\_\_\_\_\_\_\_$\UNm$；该弹簧的劲度系数 $k$ 与其自由长度 $l_{0}$（单位为 $\Um$）的关系表达式为 $k=$\_\_\_\_\_\_\_\_$\UNm$。
\end{enumerate}
\twopicture[showlabel=false, align=right, wA=3.4cm, wB=8.5cm]
{figs/23a.png}
{figs/23b.png}
%% answer:
\jdanswer{
\begin{enumerate}
	\item
	\ccnum{①} $81.7$，\ccnum{②} $0.0122$
	\item
	如图所示（$1/k-n$ 图线为过原点的直线）
	\item
	$k=\frac{1.75\times10^{3}}{n}$，$k=\frac{3.47}{l_{0}}$
\end{enumerate}
}
%% memo:
\memoanswer{
（1）$k=\frac{mg}{\Delta x_{2}}=\frac{0.1\times 9.8}{(5.26-4.06)\times 10^{-2}}\UNm=81.7\UNm$，故 $\frac{1}{k}=\frac{1}{81.7}\UNmn=0.0122\UNmn$。

（2）如图所示。

\onepicture[w=7.0cm]{figs/23_an.png}

（3）由图线可得其斜率为：$\frac{0.0347-0.0061}{60-10}=0.000572$。

故直线满足 $\frac{1}{k}=0.000572n$ 即 $k=\frac{1.75\times10^{3}}{n}$（$\UNm$）（在 $\frac{1.67\times10^{3}}{n}\sim\frac{1.83\times10^{3}}{n}$ 之间均可）。

由于 60 匝弹簧总长度为 $11.88\Ucm$；则 $n$ 匝弹簧的 $l_{0}$ 满足 $\frac{n}{l_{0}}=\frac{60}{11.88\times 10^{-2}}$。

代入 $k=\frac{1.75\times10^{3}}{n}$ 得：$k=\frac{3.47}{l_{0}}$（在 $\frac{3.31}{l_{0}}\sim\frac{3.62}{l_{0}}$ 之间均可）。
}

\item
%% number: 24
%% paperName: 2014年全国课标Ⅱ第 24 题 13 分
%% typeId: 6
%% type: 计算题
%% chapter: 运动和力
%% point: 变加速直线运动
%% method: 无
%% score: 13
%% degree: 650
%% duplicateId: 0
%% body:
2012 年 10 月，奥地利极限运动员菲利克斯·鲍姆加特纳乘气球升至约 $39\Ukm$ 的高空后跳下，经过 4 分 20 秒到达距地面约 $1.5\Ukm$ 高度处，打开降落伞并成功落地，打破了跳伞运动的多项世界纪录，取重力加速度的大小 $g=10\Umsq$。
\begin{enumerate}
	\item
	忽略空气阻力，求该运动员从静止开始下落到 $1.5\Ukm$ 高度处所需要的时间及其在此处速度的大小；
	\item
	实际上物体在空气中运动时会受到空气阻力，高速运动受阻力大小可近似表示为 $f=kv^{2}$，其中 $v$ 为速率，$k$ 为阻力系数，其数值与物体的形状，横截面积及空气密度有关。已知该运动员在某段时间内高速下落的 $v-t$ 图象如图所示，着陆过程中，运动员和所携装备的总质量 $m=100\Ukg$，试估算该运动员在达到最大速度时所受阻力的阻力系数（结果保留 1 位有效数字）。
\end{enumerate}
\onepicture[align=right, w=6.5cm]{figs/24.png}
%% answer:
\jdanswer{
\begin{enumerate}
	\item
	$87\Us$，$8.7\times10^{2}\Ums$
	\item
	$0.008\Ukgm$
\end{enumerate}
}
%% memo:
\memoanswer{
（1）设运动员从开始自由下落至 $1.5\Ukm$ 高度处的时间为 $t$，下落距离为 $h$，在 $1.5\Ukm$ 高度处的速度大小为 $v$，由运动学公式有：

$v=gt$

$h=\frac{1}{2}gt^{2}$

且 $h=3.9\times10^{4}\Um-1.5\times10^{3}\Um=3.75\times10^{4}\Um$

联立解得：$t=87\Us$，$v=8.7\times10^{2}\Ums$。

（2）运动员在达到最大速度 $v_{m}$ 时，加速度为零，由牛顿第二定律有：

$mg=kv_{m}^{2}$

由题图可读出 $v_{m}\approx 360\Ums$

代入得：$k=0.008\Ukgm$。
}

\item
%% number: 25
%% paperName: 2014年全国课标Ⅱ第 25 题 19 分
%% typeId: 6
%% type: 计算题
%% chapter: 电磁感应
%% point: 电磁感应能量分析
%% method: 无
%% score: 19
%% degree: 450
%% duplicateId: 0
%% body:
半径分别为 $r$ 和 $2r$ 的同心圆形导轨固定在同一水平面内，一长为 $r$、质量为 $m$ 且质量分布均匀的直导体棒 AB 置于圆导轨上面，BA 的延长线通过圆导轨中心 O，装置的俯视图如图所示，整个装置位于一匀强磁场中，磁感应强度的大小为 $B$，方向竖直向下，在内圆导轨的 C 点和外圆导轨的 D 点之间接有一阻值为 $R$ 的电阻（图中未画出）。直导体棒在水平外力作用下以角速度 $\omega$ 绕 O 逆时针匀速转动，在转动过程中始终与导轨保持良好接触。设导体棒与导轨之间的动摩擦因数为 $\mu$，导体棒和导轨的电阻均可忽略。重力加速度大小为 $g$。求：
\begin{enumerate}
	\item
	通过电阻 $R$ 的感应电流的方向和大小；
	\item
	外力的功率。
\end{enumerate}
\onepicture[align=right, w=4.2cm]{figs/25.png}
%% answer:
\jdanswer{
\begin{enumerate}
	\item
	C 端流向 D 端，$\frac{3\omega Br^{2}}{2R}$
	\item
	$\frac{3}{2}\mu mg\omega r+\frac{9\omega^{2}B^{2}r^{4}}{4R}$
\end{enumerate}
}
%% memo:
\memoanswer{
（1）在 $\Delta t$ 时间内，导体棒扫过的面积为：$\Delta S=\frac{1}{2}\omega\Delta t\left[(2r)^{2}-r^{2}\right]$\ccnum{①}

根据法拉第电磁感应定律，导体棒产生的感应电动势大小为：$E=\frac{B\Delta S}{\Delta t}$\ccnum{②}

根据右手定则，感应电流的方向是从 B 端流向 A 端，因此流过导体棒的电流方向是从 C 端流向 D 端；由欧姆定律流过导体棒的电流满足：$I=\frac{E}{R}$\ccnum{③}

联立\ccnum{①②③}可得：$I=\frac{3\omega Br^{2}}{2R}$\ccnum{④}。

（2）在竖直方向有：$mg-2F_{N}=0$\ccnum{⑤}

式中，由于质量分布均匀，内外圆导轨对导体棒的正压力相等，其值为 $F_{N}$，两导轨对运动的导体棒的滑动摩擦力均为：$F_{f}=\mu F_{N}$\ccnum{⑥}

在 $\Delta t$ 时间内，导体棒在内外圆导轨上扫过的弧长分别为：$l_{1}=r\omega\Delta t$\ccnum{⑦}

$l_{2}=2r\omega\Delta t$\ccnum{⑧}

克服摩擦力做的总功为：$W_{f}=F_{f}(l_{1}+l_{2})$\ccnum{⑨}

在 $\Delta t$ 时间内，消耗在电阻 $R$ 上的功为：$W_{R}=I^{2}R\Delta t$\ccnum{⑩}

根据能量转化和守恒定律，外力在 $\Delta t$ 时间内做的功为 $W=W_{f}+W_{R}$\ccnum{⑪}

外力的功率为：$P=\frac{W}{\Delta t}$\ccnum{⑫}

由\ccnum{④}至\ccnum{⑫}式可得：$P=\frac{3}{2}\mu mg\omega r+\frac{9\omega^{2}B^{2}r^{4}}{4R}$\ccnum{⑬}。
}

\item
%% number: 33
%% paperName: 2014年全国课标Ⅱ第 33 题 10 分
%% typeId: 6
%% type: 计算题
%% chapter: 气体性质
%% point: 气体多过程
%% method: 无
%% score: 10
%% degree: 450
%% duplicateId: 0
%% body:
\textbf{（1）}（5 分）下列说法正确的是\xzanswer{BCE}
\fivechoices[answer=BCE]
{悬浮在水中的花粉的布朗运动反映了花粉分子的热运动}
{空中的小雨滴呈球形是水的表面张力作用的结果}
{彩色液晶显示器利用了液晶的光学性质具有各向异性的特点}
{高原地区水的沸点较低，这是高原地区温度较低的缘故}
{干湿泡温度计的湿泡显示的温度低于干泡显示的温度，这是湿泡外纱布中水蒸发吸热的结果}

\textbf{（2）}（10 分）如图，两气缸 A、B 粗细均匀、等高且内壁光滑。其下部由体积可忽略的细管连通；A 的直径是 B 的 2 倍，A 上端封闭，B 上端与大气连通；两气缸除 A 顶部导热外，其余部分均绝热。两气缸中各有一厚度可忽略的绝热轻活塞 a、b，活塞下方充有氮气，活塞 a 上方充有氧气。当大气压为 $p_{0}$，外界和气缸内气体温度均为 $7\UdC$ 且平衡时，活塞 a 离气缸顶的距离是气缸高度的 $\frac{1}{4}$，活塞 b 在气缸正中间。
\begin{enumerate}
	\item
	现通过电阻丝缓慢加热氮气，当活塞 b 恰好升至顶部时，求氮气的温度；
	\item
	继续缓慢加热，使活塞 a 上升，当活塞 a 上升的距离是气缸高度的 $\frac{1}{16}$ 时，求氧气的压强。
\end{enumerate}
\onepicture[align=right, w=4.5cm]{figs/33.png}
%% answer:
\jdanswer{
\begin{enumerate}
	\item
	$320\UK$
	\item
	$\frac{4}{3}p_{0}$
\end{enumerate}
}
%% memo:
\memoanswer{
（1）悬浮在水中的花粉的布朗运动反映了水分子的无规则热运动，A 错误；空中的小雨滴呈球形是水的表面张力作用的结果，B 正确；彩色液晶显示器利用了液晶的光学性质具有各向异性的特点，C 正确；高原地区水的沸点较低，是由于高原地区气压低，故水的沸点也较低，D 错误；干湿泡温度计的湿泡显示的温度低于干泡显示的温度，是由于湿泡外纱布中水蒸发吸收热量，使湿泡的温度降低的缘故，E 正确。

（2）（i）活塞 b 升至顶部的过程中，活塞 a 不动，活塞 a、b 下方的氮气经历等压过程，设气缸 A 的容积为 $V_{0}$，氮气初始状态的体积为 $V_{1}$，温度为 $T_{1}$，末态体积为 $V_{2}$，温度为 $T_{2}$，按题意，气缸 B 的容积为 $V_{0}/4$，由题给数据及盖吕萨克定律有：$\frac{V_{1}}{T_{1}}=\frac{V_{2}}{T_{2}}$\ccnum{①}

且 $V_{1}=\frac{3}{4}V_{0}+\frac{1}{2}\cdot\frac{V_{0}}{4}=\frac{7}{8}V_{0}$\ccnum{②}

$V_{2}=\frac{3}{4}V_{0}+\frac{V_{0}}{4}=V_{0}$\ccnum{③}

由\ccnum{①②③}式及所给的数据可得：$T_{2}=320\UK$\ccnum{④}。

（ii）活塞 b 升至顶部后，由于继续缓慢加热，活塞 a 开始向上移动，直至活塞上升的距离是气缸高度的 $\frac{1}{16}$ 时，活塞 a 上方的氧气经历等温过程，设氧气初始状态的体积为 $V_{1}'$，压强为 $p_{1}'$；末态体积为 $V_{2}'$，压强为 $p_{2}'$，由所给数据及玻意耳定律可得

$V_{1}'=\frac{1}{4}V_{0}$，$p_{1}'=p_{0}$，$V_{2}'=\frac{3}{16}V_{0}$\ccnum{⑤}

$p_{1}'V_{1}'=p_{2}'V_{2}'$\ccnum{⑥}

由\ccnum{⑤⑥}式可得：$p_{2}'=\frac{4}{3}p_{0}$\ccnum{⑦}。
}

\item
%% number: 34
%% paperName: 2014年全国课标Ⅱ第 34 题 10 分
%% typeId: 6
%% type: 计算题
%% chapter: 光学
%% point: 光的折射
%% method: 无
%% score: 10
%% degree: 450
%% duplicateId: 0
%% body:
\textbf{（1）}（5 分）图（a）为一列简谐横波在 $t=0.10\Us$ 时的波形图，P 是平衡位置在 $x=1.0\Um$ 处的质点，Q 是平衡位置在 $x=4.0\Um$ 处的质点；图（b）为质点 Q 的振动图形。下列说法正确的是\xzanswer{BCE}
\onepicture[w=12.0cm]{figs/34.png}
\fivechoices[answer=BCE]
{在 $t=0.10\Us$ 时，质点 Q 向 $y$ 轴正方向运动}
{在 $t=0.25\Us$ 时，质点 P 的加速度方向与 $y$ 轴正方向相同}
{从 $t=0.10\Us$ 到 $t=0.25\Us$，该波沿 $x$ 轴负方向传播了 $6\Um$}
{从 $t=0.10\Us$ 到 $t=0.25\Us$，质点 P 通过的路程为 $30\Ucm$}
{质点 Q 简谐运动的表达式为 $y=0.10\sin10\pi t$（国际单位制）}

\textbf{（2）}（10 分）一厚度为 $h$ 的大平板玻璃水平放置，下表面贴有一半径为 $r$ 的圆形发光面。在玻璃板上表面放置一半径为 $R$ 的圆纸片，圆纸片与圆形发光面的中心在同一竖直线上。已知圆纸片恰好能完全遮挡住从圆形发光面发出的光线（不考虑反射），求平板玻璃的折射率。
%% answer:
\jdanswer{
$\sqrt{1+\left(\frac{h}{R-r}\right)^{2}}$
}
%% memo:
\memoanswer{
（1）由 Q 点的振动图线可知，$t=0.10\Us$ 时质点 Q 向 $y$ 轴负方向振动，A 错误；由波的图像可知，波向左传播，波的周期为 $T=0.2\Us$，$t=0.10\Us$ 时质点 P 向上振动，经过 $0.15\Us=3T/4$ 时，即在 $t=0.25\Us$ 时，质点振动到 $x$ 轴下方位置，且速度方向向上，加速度方向也沿 $y$ 轴正向，B 正确；波速 $v=\frac{\lambda}{T}=\frac{8}{0.2}\Ums=40\Ums$，故从 $t=0.10\Us$ 到 $t=0.25\Us$，该波沿 $x$ 负方向传播的距离为：$x=vt=40\times0.15\Um=6\Um$，C 正确；由于 P 点不是在波峰或波谷或者平衡位置，故从 $t=0.10\Us$ 到 $t=0.25\Us$ 的 $3/4$ 周期内，通过的路程不等于 $3A=30\Ucm$，选项 D 错误；质点 Q 做简谐振动的表达式为：$y=A\sin\left(\frac{2\pi}{T}\right)t=0.10\sin10\pi t$（国际单位），选项 E 正确。

（2）如图，考虑从圆形发光面边缘的 A 点发出的一条光线，假设它斜射到玻璃上表面的 A 点折射，根据折射定律有：$n\sin\theta=\sin\alpha$

式中，$n$ 是玻璃的折射率，$\theta$ 是入射角，$\alpha$ 是折射角。

现假设 A 恰好在纸片边缘，由题意，在 A 刚好发生全反射，故 $\alpha=\frac{\pi}{2}$。

设 AA' 线段在玻璃上表面的投影长为 $L$，由几何关系有：$\sin\theta=\frac{L}{\sqrt{L^{2}+h^{2}}}$

由题意纸片的半径应为 $R=L-r$

联立以上各式可得：$n=\sqrt{1+\left(\frac{h}{R-r}\right)^{2}}$。

\onepicture[w=5.0cm]{figs/34_an.png}
}

\item
%% number: 35
%% paperName: 2014年全国课标Ⅱ第 35 题 10 分
%% typeId: 4
%% type: 实验题
%% chapter: 运动和力
%% point: 动量守恒定律
%% method: 无
%% score: 10
%% degree: 450
%% duplicateId: 0
%% body:
\textbf{（1）}（5 分）在人类对微观世界进行探索的过程中，科学实验起到了非常重要的作用。下列说法符合历史事实的是\xzanswer{ACE}
\fivechoices[answer=ACE]
{密立根通过油滴实验测出了基本电荷的数值}
{贝克勒尔通过对天然放射现象的研究，发现了原子中存在原子核}
{居里夫妇从沥青铀矿中分离出钋（Po）和镭（Ra）两种新元素}
{卢瑟福通过 $\alpha$ 粒子散射实验证实了原子核内部存在质子}
{汤姆逊通过阴极射线在电场和磁场中的偏转实验，发现了阴极射线是由带负电的粒子组成的，并测出了该粒子的比荷}

\textbf{（2）}（10 分）现利用图（a）所示装置验证动量守恒定律。在图（a）中，气垫导轨上有 A、B 两个滑块，滑块 A 右侧带有一弹簧片，左侧与打点计时器（图中未画出）的纸带相连；滑块 B 左侧也带有一弹簧片，上面固定一遮光片，光电计数器（未完全画出）可以记录遮光片通过光电门的时间。

实验测得滑块 A 的质量 $m_{1}=0.301\Ukg$，滑块 B 的质量 $m_{2}=0.108\Ukg$，遮光片的宽度 $d=1.00\Ucm$；打点计时器所用交流电的频率 $f=50.0\UHz$。

将光电门固定在滑块 B 的右侧，启动打点计时器，给滑块 A 一向右的初速度，使它与 B 相碰。碰后光电计数器显示的时间为 $\Delta t_{B}=3.500\UmsT$，碰撞前后打出的纸带如图（b）所示。

若实验允许的相对误差绝对值（$\left|\frac{\text{碰撞前后动量之差}}{\text{碰前总动量}}\right|\times100\%$）最大为 $5\%$，本实验是否在误差范围内验证了动量守恒定律？写出运算过程。
\twopicture[showlabel=false, align=right, wA=7.5cm, wB=12.0cm]
{figs/35a.png}
{figs/35b.png}
%% answer:
\jdanswer{
在允许的误差范围内验证了动量守恒定律，$\delta_{r}=1.7\%<5\%$。
}
%% memo:
\memoanswer{
（1）密立根通过油滴实验测出了基本电荷的数值为 $1.6\times10^{-19}\UC$，A 正确；贝克勒尔发现了天然放射现象，而不是原子中存在原子核，B 错误；居里夫妇从沥青铀矿中分离出了钋（Po）和镭（Ra）两种新元素，C 正确；卢瑟福通过 $\alpha$ 粒子散射实验，得出了原子的核式结构理论，D 错误；汤姆逊通过对阴极射线在电场及在磁场中偏转的实验，发现了阴极射线是由带负电的粒子组成，并测定了粒子的比荷，E 正确。

（2）按定义，物体运动的瞬时速度大小 $v$ 为：$v=\frac{\Delta x}{\Delta t}$；式中 $\Delta x$ 为物块在很短的时间 $\Delta t$ 内的位移，设纸带上打出相邻两点的时间间隔为 $\Delta t_{A}$，则 $\Delta t_{A}=1/f=0.02\Us$，可视为很短。设在 A 碰撞前后瞬时速度大小分别为 $v_{0}$ 和 $v_{1}$，由图（b）所给数据可得 $v_{0}=2.00\Ums$，$v_{1}=0.790\Ums$。设 B 碰撞后瞬时速度大小为 $v_{2}$，$v_{2}=\frac{d}{\Delta t_{B}}=2.56\Ums$。设两滑块在碰撞前后的动量分别为 $P$ 和 $P'$，则 $P=m_{1}v_{0}$，$P'=m_{1}v_{1}+m_{2}v_{2}$。两滑块在碰撞前后总动量相对误差的绝对值为 $\delta_{r}=\left|\frac{P-P'}{P}\right|\times100\%$，联立各式代入数据得 $\delta_{r}=1.7\%<5\%$，因此，本实验在允许的误差范围内验证了动量守恒定律。
}

\end{enumerate}

\end{document}
